\documentclass[a4paper]{article}
\usepackage[margin=1.5cm,top=1.5cm,bottom=1.5cm,a4paper]{geometry}
\usepackage{graphicx} % Required for inserting images
\usepackage{titlesec} % Required for customizing section titles
\usepackage[margin=1.5cm,top=1.5cm, bottom=1.5cm]{geometry} % Ajusta los márgenes aquí
\usepackage{multicol} % Required for multicols environment
\usepackage{parskip}
\usepackage{etoolbox}
\usepackage{bm}
\usepackage{tcolorbox}
\usepackage{float}
\usepackage{amsmath}
\usepackage{dingbat}
\usepackage{booktabs} 
\usepackage{caption}
\usepackage[utf8]{inputenc}
\usepackage[T1]{fontenc}
\usepackage{amssymb}
\usepackage[spanish]{babel}
\usepackage{listings}
\usepackage{xcolor}

\usepackage[backend=biber,style=ieee,citestyle=ieee]{biblatex}

\addbibresource{referencias.bib}


\lstset{
	language=Matlab,
	basicstyle=\ttfamily\small,
	keywordstyle=\color{blue},
	commentstyle=\color{gray},
	numbers=left,
	numberstyle=\tiny,
	frame=single
}

\addto\captionsspanish{
	\renewcommand{\tablename}{Tabla}
}

\usepackage{tikz}
\usetikzlibrary{arrows.meta, positioning}
\usetikzlibrary{babel}

\usepackage{csquotes}






% Personalización de la sección para que comience desde la izquierda
\titleformat{\section}[hang]{\normalfont\Large\bfseries}{\thesection}{1em}{}

\begin{document}
	
	% ---------------- PORTADA ----------------
	\begin{titlepage}
		\centering
		
		% Logo institucional
		\vspace*{0.8cm}
		\includegraphics[height=1.5cm]{img/logofiuba.png}
		
		\vspace{1cm}
		
		% Información de la materia
		{\Large \textsc{(TA132) Procesos Estocásticos}\par}
		
		\vspace{0.8cm}
		
		{\large Trabajo Práctico\par}
		
		\vspace{0.5cm}
		
		% Título principal
		\rule{0.85\textwidth}{0.6pt}
		
		\vspace{0.7cm}
		
		{\Huge\bfseries Taller de estimación\par}
		
		\vspace{0.7cm}
		
		\rule{0.85\textwidth}{0.6pt}
		
		\vspace{2cm}
		
		% Autores
		{\Large\bfseries Autores\par}
		
		\vspace{0.7cm}
		
		{\large
			106213  Sebastián Lazo\\
			\texttt{slazo@fi.uba.ar}\\[0.6cm]
			Kevin Jesús Sterling Pérez\\
			\texttt{-----}
			\par}
		
		\vfill
		
		% Información institucional
		\rule{0.6\textwidth}{0.4pt}
		
		\vspace{0.5cm}
		
		{\large Facultad de Ingeniería\\
			Universidad de Buenos Aires\par}
		
		\vspace{0.8cm}
		
		{\large 23 de noviembre de 2026\par}
		
		\vspace{0.5cm}
		
	\end{titlepage}
 \newpage
 
 \section*{Introducción}
 
 El informe describe el desarrollo de un taller dentro del ámbito de la materia Procesos estocásticos en el cual se abarca a nivel teórico y práctico la metodología de estimación de procesos aleatorios. Se encuentra dividido en tres actividades principales, en la primera actividad, utilizando bibliografía especifica, se estudian, simulan y analizan diferentes ejemplos, la segunda actividad se utiliza un \textit{paper} acerca de la modulación OFDM se relaciona con lo previo y se propone un experimento numérico que permita evaluar la bondad de dichas aproximaciones. La actividad tres finalmente es el desarrollo del informe.
 
 \subsection*{Contexto}
 
Un proceso aleatorio (o estocástico) es una familia de variables aleatorias ${X(t, \omega) : t \in T}$ definida sobre un mismo espacio de probabilidad $(\Omega, F, P)$, indexada por un parámetro t en un conjunto T (usualmente el tiempo) o n perteneciente al conjunto $\mathbb{Z}$ describiendo usualmente tiempo discreto:


\begin{equation}
	 \{X(t, \omega) : t \in T,\, \omega \in \Omega\}, \qquad X : T \times \Omega \to \mathbb{R}^d
\end{equation}
 
 
\begin{equation}
	\{X[n, \omega] : n \in \mathbb{Z},\, \omega \in \Omega\}, \qquad X : \mathbb{Z} \times \Omega \to \mathbb{R}^d
\end{equation}


Un proceso estocástico estacionario en sentido amplio (ESA) es un proceso aleatorio $\{X[n]\}$ si se cumple simultáneamente:

\begin{equation}
	\mu_X[n] = \mathbb{E}\{X[n]\} = \mu_X, \qquad \forall n \in \mathbb{Z}
\end{equation}

\begin{equation}
	R_{XX}[n_1, n_2] = \mathbb{E}\{X[n_1]X^*[n_2]\} = R_{XX}[n_1 - n_2] = R_{XX}[k]
\end{equation}

Siendo:
\begin{equation}
	\mu_X[n] = \mathbb{E}\{X[n]\} = \int_{-\infty}^{\infty} x \, f_{X[n]}(x) \, dx
\end{equation}

\begin{equation}
	R_{XX}[n_1, n_2] = \mathbb{E}\{X[n_1] X^*[n_2]\} = \iint_{-\infty}^{\infty} x_1 x_2^* \, f_{X[n_1], X[n_2]}(x_1, x_2) \, dx_1 \, dx_2
\end{equation}



\subsection*{Estimación}

Dados dos procesos estocásticos $X(n)$ e $Y(n)$ ESA. La estimación consiste en inferir realización desconocidas de $X(n)$ conociendo un conjunto de realizaciones de $\textbf{Y}(n)$ de nomenclatura $\vec{y}(n)$ durante una ventana de tiempo $n \in I$. El estimador es la función $g(\vec{y}(n))$, preferentemente lineal, que relaciona estas realizaciones de $Y(n)$ con los valores estimados de $X(n)$. La fidelidad de la estimación se obtiene conociendo el error cuadrático mínimo (MSE \textit{mean squared error} )y buscando la función $g(\vec{y}(n))$ que minimice el error (MMSE \textit{minimum mean squared error} ), se puede demostrar que la estimación lineal que cumple esto es: 

\begin{equation}
	\hat{x}(n) = [R_Y^{-1} \cdot r_{XY}]^{T} \cdot \vec{y}(n)
\end{equation}

Siendo: 

\begin{equation}
	r_{XY} = \mathbb{E}[X(n)\cdot \textbf{Y}(n)]  \quad \wedge \quad R_{Y} = \mathbb{E}[\textbf{Y}(n)\cdot \textbf{Y}^T(n)]
\end{equation}


 \section*{Actividad 1: Análisis y simulación de sistemas de filtrado óptimo }
 Se utiliza la bibliografía \cite{hayes1996} en donde se menciona la estimación en el ámbito de las señales, tomando cada proceso estocástico como un conjunto de señales relacionadas con variables aleatorias. La estimación de una señal teniendo valores de otra es uno de los mas importantes desafíos en el procesamiento de señales y abarca interesantes aplicaciones.
 
 Se toma como ejemplo la recuperación de una señal que se envío a través de un medio no ideal en una comunicación digital. En situaciones simples, es posible utilizar filtros clásicos, como los pasabajos, pasaaltos o pasabandas, para recuperar la señal. Sin embargo, estos filtros no necesariamente proporcionan la mejor estimación posible. Por este motivo, la bibliografía introduce el diseño de filtros digitales óptimos, entre ellos el filtro de Wiener y el filtro de Kalman discreto.
 
 \subsection*{Filtros Wiener}
 
 Norbert Wiener investiga el problema del diseño de filtros para producir una estimación óptima de la señal desde una observación con ruido. Se define $x(n)$ como la observación con ruido, $d(n)$ la señal a recuperar, $\hat{d}(n)$ la estimación obtenida y $v(n)$ el ruido y $e(n) = d(n) - \hat{d}(n)$ como el error en la estimación.
 
 \begin{equation}
 	x(n) = d(n) + v(n)
 \end{equation}

 Se asume $d(n)$ y $v(n)$ como procesos ESA, Wiener considera el problema de diseño de la misma forma que se menciona en la introducción del informe, buscando minimizar el error cuadrático medio mediante un filtro $W(z)$.

\begin{figure}[H]
	\begin{tikzpicture}[transform shape]
		% Paths, nodes and wires:
		\node[shape=rectangle, draw, line width=1pt, minimum width=4.215cm, minimum height=2.215cm] at (27.375, 24.125){} node[anchor=center, align=center, text width=3.777cm, inner sep=6.7pt] at (27.375, 24.125){\Huge $W(z)$};
		\node[shape=rectangle, minimum width=2.215cm, minimum height=1.215cm] at (21.625, 24.25){} node[anchor=north west, align=left, text width=1.827cm, inner sep=6pt] at (20.5, 24.875){\Huge $x(n)$\\};
		\draw[line width=2pt, -latex] (22.75, 24.125) |- (25.25, 24.125);
		\draw[line width=2pt, -latex] (29.5, 24.125) |- (32, 24.125);
		\node[shape=rectangle, minimum width=2.215cm, minimum height=1.215cm] at (30.625, 24.875){} node[anchor=north west, align=left, text width=1.827cm, inner sep=6pt] at (29.5, 25.5){\huge $\hat{d}(n)$\\};
		\node[shape=rectangle, minimum width=2.215cm, minimum height=1.215cm] at (32.75, 27.625){} node[anchor=north west, align=left, text width=1.827cm, inner sep=6pt] at (31.625, 28.25){\huge $d(n)$\\};
		\draw[line width=2pt, -latex] (32.49, 27.115) -| (32.49, 24.615);
		\draw[line width=2pt, -latex] (32.98, 24.125) |- (35.48, 24.125);
		\node[shape=rectangle, minimum width=2.215cm, minimum height=1.215cm] at (36.875, 24.125){} node[anchor=north west, align=left, text width=1.827cm, inner sep=6pt] at (35.75, 24.75){\huge $e(n)$\\};
		\node[shape=rectangle, minimum width=0.965cm, minimum height=0.465cm] at (31.75, 23.75){} node[anchor=north west, align=left, text width=0.577cm, inner sep=6pt] at (31.25, 24){\Huge -};
		\node[shape=rectangle, minimum width=0.965cm, minimum height=0.465cm] at (33.25, 24.75){} node[anchor=north west, align=left, text width=0.668cm, inner sep=4.7pt] at (32.75, 25){\Large +};
		\node[shape=circle, draw, line width=1pt, minimum width=0.715cm](N1) at (32.5, 24.125){} node[anchor=center] at (N1.text){$+$};
	\end{tikzpicture}
\end{figure}


Se debe diseñar un filtro lineal invariante en el tiempo LTI, $W(z)$, que filtre una señal dada, $x(n)$, para producir la estimación mínima de la media cuadrática, $\hat{d}(n)$, de otra señal, $d(n)$. Dependiendo de cómo se relacionen las señales $x(n)$ y $d(n)$, se pueden plantear varios problemas diferentes e importantes dentro de un marco de filtrado de Wiener. Se considerarán en la bibliografía:

\begin{itemize}
	\item \textbf{Filtrado:} Este es el problema clásico de la introducción, en el que se nos da $x(n) = d(n) + v(n)$ y el objetivo es estimar $d(n)$ utilizando un filtro causal, es decir, estimar $d(n)$ a partir de los valores actuales y pasados de $x(n)$.
	
	\item \textbf{Suavizado (\textit{smoothing}):} Esto es lo mismo que el problema de filtrado, excepto que se permite que el filtro no sea causal. Un filtro de suavizado de Wiener, por ejemplo, puede diseñarse para estimar $d(n)$ a partir de $x(n) = d(n) + v(n)$ utilizando todos los datos disponibles.
	
	\item \textbf{Predicción:} Si $d(n) = x(n + 1)$ y $W(z)$ es un filtro causal, entonces el filtro de Wiener se convierte en un predictor lineal. En este caso, el filtro debe producir una predicción (estimación) de $x(n + 1)$ en términos de una combinación lineal de valores anteriores de $x(n)$.
	
	\item \textbf{Deconvolución:} Cuando $x(n) = d(n) \ast g(n) + v(n)$, donde $g(n)$ es la respuesta de muestra unitaria de un filtro LTI, el filtro de Wiener se convierte en un filtro de deconvolución.
\end{itemize}


\subsubsection*{7.2 : Filtros Wiener de respuesta al impulso finita (FIR)}

Se considera el tipo de filtros de Wiener FIR $W(z)$, obteniendo el MMSE, para esta sección se asumen nuevamente $x(n)$ y $d(n)$ procesos ESA y con funciones de autocorrelación conocidas $r_x(k)$ , $r_{d}(k)$ y la función de auto-correlación cruzada $r_{dx}(k)$.

El error mínimo cuadrado es definido como:

\begin{equation}
	MSE = \mathbb{E}[|d(n)- \hat{d}(n)|^2] =  \mathbb{E}[|e(n)|^2]
\end{equation}

Donde se obtiene la estimación por medio de la salida del filtro Wiener $\hat{d}(n)$ = $w(n) \ast x(n) $. 
 
Se puede deducir que para minimizar el error se debe encontrar y pensando el error como una función de $MSE(w[k])$, siendo $w[k]$ el $k$-ésimo coeficiente del filtro un numero complejo y $w^*[k]$ su conjugado, los puntos estacionarios de dicha función, los mínimos locales y globales, deben satisfacer la siguiente condición: 
 
 \begin{equation}
 		\frac{\partial MSE(w[k])}{\partial w^*[k]} = 0 
 \end{equation}

 Utilizando: 
 \begin{equation}
 	|e(n)|^{2} = e(n)\cdot e^{\ast} (n) 
 \end{equation}


  
 \begin{equation}
 	\frac{\partial MSE(w[k])}{\partial w^*[k]} = \frac{\partial }{\partial w^*[k]}  \mathbb{E}[|e(n)|^2] =\frac{\partial }{\partial w^*[k]}  \mathbb{E}[e(n)\cdot e^{\ast} (n) ]  = 0 
 \end{equation}
	
  Con 
  \begin{equation}
  	e(n) = d(n) - \sum_{m=0}^{p} w(m)\,x(n-m)
  \end{equation}
  
  \begin{equation}
  	e^*(n) = d^*(n) - \sum_{m=0}^{p} w^*(m)\,x^*(n-m)
  \end{equation}
	

	
  Se desplaza la derivada dentro de la esperanza y afecta solo a $	e^*(n) $:

  
	\begin{equation}
		\mathbb{E}\Bigl[e(n)\cdot\frac{\partial e^{\ast}(n)}{\partial w^*[k]}\Bigr] = 0
	\end{equation}

Desarrollando la derivada:

\begin{equation}
	\frac{\partial e^{\ast}(n)}{\partial w^*[k]} = \frac{	\partial }{\partial w^*[k]}[d^*(n) - \sum_{m=0}^{p} w^*(m)\,x^*(n-m)]
\end{equation}

Donde en la sumatoria sobrevive solo el termino $w^*(m) =w^*(k)$

\begin{equation}
	\frac{\partial e^{\ast}(n)}{\partial w^*[k]} = \frac{	\partial }{\partial w^*[k]}[d^*(n) -  w^*(k)\,x^*(n-m)]
\end{equation}

Derivando se descarta $d^{\ast}(n)$ obteniendo:

\begin{equation}
	\frac{\partial e^{\ast}(n)}{\partial w^*[k]} = - x^*(n-m)
\end{equation}

Volviendo a introducirlo dentro de la esperanza quitando el signo al multiplicar ambos lados por -1:

\begin{equation}
	\mathbb{E}\Bigl \{ e(n)\cdot [x^*(n-m)]\Bigr \} = 0
\end{equation}

Utilizando la descripción de $e(n)$ previa:

\begin{equation}
	\mathbb{E}\Bigl \{ \Bigl (d(n) - \sum_{m=0}^{p} w(m)\,x(n-m) \Bigr )\cdot [x^*(n-m)] \Bigr \} = 0
\end{equation}


\begin{equation}
	\mathbb{E}\Bigl \{  d(n)[x^*(n-m)] - \sum_{m=0}^{p} w(m)\,x(n-m)[x^*(n-m)] \Bigr \} = 0
\end{equation}




\begin{equation}
	\mathbb{E}\Bigl \{  d(n)[x^*(n-m)]\Bigr \} - \mathbb{E} \Bigl \{ \sum_{m=0}^{p} w(m)\,x(n-m)[x^*(n-m)] \Bigr \} = 0
\end{equation}



\begin{equation}
	\mathbb{E}\Bigl \{  d(n)[x^*(n-m)]\Bigr \} -  \sum_{m=0}^{p} w(m)\, \mathbb{E} \Bigl \{ x(n-m)[x^*(n-m)] \Bigr \} = 0
\end{equation}


Finalmente, dado que $x(n)$ y $d(n)$ son conjuntamente ESA, entonces $\mathbb{E} \{x(n- m )x*(n-k)\} =r_x(k-m)$ y $\mathbb{E} \{d(n)x*(n-k)\} = r_{dx}(k)$ y la ecuación se convierte en : 



\begin{equation}
  r_{dx}(k)  -  \sum_{m=0}^{p} w(m)\, r_x(k-m) \Bigr \} = 0
\end{equation}


\begin{equation}
  \sum_{m=0}^{p} w(m)\, r_x(k-m)  =  r_{dx}(k)
\end{equation}
que es un conjunto de $p$ ecuaciones lineales en las $p$ incógnitas $w(k)$, $k = 0, 1, \dots, p-1$. En forma matricial, utilizando la simetría conjugada de la secuencia de autocorrelación, $r_x(-k) = r_x^*(k)$, la ecuación se convierte en un sistema lineal matricial:

\begin{equation}
	\begin{bmatrix}
		r_x(0)     & r_x^*(1)   & \cdots & r_x^*(p-1) \\
		r_x(1)     & r_x(0)     & \cdots & r_x^*(p-2) \\
		r_x(2)     & r_x(1)     & \cdots & r_x^*(p-3) \\
		\vdots     & \vdots     &        & \vdots     \\
		r_x(p-1)   & r_x(p-2)   & \cdots & r_x(0)
	\end{bmatrix}
	\begin{bmatrix}
		w(0) \\
		w(1) \\
		w(2) \\
		\vdots \\
		w(p-1)
	\end{bmatrix}
	=
	\begin{bmatrix}
		r_{dx}(0) \\
		r_{dx}(1) \\
		r_{dx}(2) \\
		\vdots \\
		r_{dx}(p-1)
	\end{bmatrix}
\end{equation}


Comprimiendo los vectores y la matriz con la notación $\mathbf{R}_x$ , $\mathbf{w} $ y $ \mathbf{r}_{dx}$:

\begin{equation}
	\boxed{\mathbf{R}_x \mathbf{w} = \mathbf{r}_{dx}}
\end{equation}

De esta expresión se puede despejar los coeficientes del vector $\mathbf{w} $:

\begin{equation}
	\boxed{ \mathbf{w} = \mathbf{R}_x^{-1}\mathbf{r}_{dx}}
\end{equation}

Llegando así a la misma expresión presentada en la introducción. Para obtener las estimaciones: 


\begin{equation}
	\boxed{ \hat{d}(n) = \mathbf{w}^T \cdot x(n) }
\end{equation}

De esta forma en el libro se presentan en la siguiente tabla las ecuaciones conocidas como ecuaciones de Wiener-Hopf y las relaciones obtenidas.


\begin{table}[H]
	\centering
	\caption{The Wiener-Hopf Equations for the FIR Wiener Filter and the Minimum Mean-Square Error}
	\begin{tabular}{ll}
		\toprule
		\textit{Wiener-Hopf equations} & \\
		& $\displaystyle \sum_{l=0}^{p-1} w(l)\,r_x(k - l) = r_{dx}(k) \quad ; \quad k = 0, 1, \dots, p-1$ \\[12pt]
		\textit{Correlations} & \\
		& $r_x(k) = E\{x(n)\,x^*(n - k)\}$ \\
		& $r_{dx}(k) = E\{d(n)\,x^*(n - k)\}$ \\[10pt]
		\textit{Minimum error} & \\
		& $\displaystyle \xi_{\min} = r_d(0) - \sum_{l=0}^{p-1} w(l)\,r_{dx}^*(l)$ \\[10pt]
		\bottomrule
	\end{tabular}
\end{table}

\begin{equation}
	\boxed{MMSE  = r_d(0) - \mathbf{r}_{dx}^H \mathbf{R}_x^{-1} \mathbf{r}_{dx}}
\end{equation}


Observando los resultados previos, se pueden obtener los valores de $\mathbf{w}$ mediante el ruido $v(n)$ presentado en la introducción ($x(n) = d(n) + v(n)$). Se obtienen la relación de sus funciones de autocorrelación.


\begin{align}
	r_{dx}(k) &= E\{d(n)x^*(n-k)\} \nonumber \\
	&= E\{d(n)d^*(n-k)\} + E\{d(n)v^*(n-k)\} \nonumber \\
	&= r_d(k)
\end{align}

A continuación, dado que:

\begin{align}
	r_x(k) &= E\{x(n+k)x^*(n)\} \nonumber \\
	&= E\bigl\{[d(n+k) + v(n+k)][d(n) + v(n)]^*\bigr\}
\end{align}

\begin{equation}
	r_x(k) = r_d(k) + r_v(k)
\end{equation}

Formando las matrices con todos los valores de k, se obtiene :

\begin{equation}
	\boxed{ r_d = [\mathbf{R}_d +  \mathbf{R}_v] \cdot \mathbf{w} }
\end{equation}



 \subsubsection*{7.2.1 : Ejemplo de filtro y simulación}
 
 
 Sea $d(n)$ un proceso $\mathrm{AR}$ con secuencia de autocorrelación
 
 
 \begin{equation}
 	r_d(k) = \alpha^{|k|}
 \end{equation}
 con $0 < \alpha < 1$, y suponiendo que $d(n)$ se observa en presencia de ruido blanco no correlacionado, $v(n)$, que tiene una varianza $\sigma_v^2$:
 \begin{equation}
 	x(n) = d(n) + v(n)
 \end{equation}
 
 Diseñemos un filtro de Wiener FIR de primer orden para reducir el ruido en $x(n)$. Con
 \begin{equation}
 	W(z) = w(0) + w(1)z^{-1}
 \end{equation}
 las ecuaciones de Wiener-Hopf son
 \begin{equation}
 	\begin{bmatrix}
 		r_x(0) & r_x(1) \\
 		r_x(1) & r_x(0)
 	\end{bmatrix}
 	\begin{bmatrix}
 		w(0) \\
 		w(1)
 	\end{bmatrix}
 	=
 	\begin{bmatrix}
 		r_{dx}(0) \\
 		r_{dx}(1)
 	\end{bmatrix}
 \end{equation}
 
 Dado que se supone que $d(n)$ y $v(n)$ no están correlacionados, entonces
 \begin{equation}
 	r_{dx}(k) = r_d(k) = \alpha^{|k|}
 \end{equation}
 y
 \begin{equation}
 	r_x(k) = r_d(k) + r_v(k) = \alpha^{|k|} + \sigma_v^2 \delta(k)
 \end{equation}
 
 Por lo tanto, las ecuaciones de Wiener-Hopf se convierten en[
 
 \begin{equation}
 	\begin{bmatrix}
 		1 + \sigma_v^2 & \alpha \\
 		\alpha & 1 + \sigma_v^2
 	\end{bmatrix}
 	\begin{bmatrix}
 		w(0) \\
 		w(1)
 	\end{bmatrix}
 	=
 	\begin{bmatrix}
 		1 \\
 		\alpha
 	\end{bmatrix}
 \end{equation}




\subsubsection*{7.2.2 : Predicción lineal}

 \subsubsection*{7.2.6 : Cancelación de ruido y simulación}
 
 
 \section*{Actividad 2}
 
 
 
 
 \section*{Conclusión}
 
 
\printbibliography
 
\end{document}